% BHU PYQ -- Professional Exam Template
\documentclass[11pt,a4paper]{article}

\usepackage[a4paper,top=2.5cm,bottom=2.5cm,left=2.8cm,right=2.8cm,headheight=14pt]{geometry}
\usepackage{amsmath,amssymb}
\usepackage[T1]{fontenc}
\usepackage[utf8]{inputenc}
\usepackage{lmodern}
\usepackage{microtype}
\usepackage{fancyhdr}
\usepackage{enumitem}
\usepackage{hyperref}
\usepackage{lastpage}
\usepackage{parskip}

\hypersetup{colorlinks=true,urlcolor=black,linkcolor=black,
  pdftitle={B.Sc. (Hons.) Botany Sem-II BOB-203A 2013-14 -- BHU PYQ}}

\pagestyle{fancy}
\fancyhf{}
\renewcommand{\headrulewidth}{0.4pt}
\renewcommand{\footrulewidth}{0.4pt}
\fancyhead[L]{\small   \textit{Banaras Hindu University}}
\fancyhead[R]{\small   \textit{Botany Paper: BOB-203A (2013-14)}}
\fancyfoot[C]{\small Page  \thepage\ of \pageref{LastPage}}


\newlist{parts}{enumerate}{2}
\setlist[parts,1]{label=(\alph*),leftmargin=*,itemsep=4pt,topsep=3pt}
\setlist[parts,2]{label=(\roman*),leftmargin=*,itemsep=2pt,topsep=2pt}

\newcommand{\pts}[1]{\hfill   \textit{[#1]}}


\begin{document}


\begin{center}
  {\large\bfseries BANARAS HINDU UNIVERSITY}\\[2pt]
  {
\normalsize B.Sc. (Hons.) Semester II Examination, 2013-14}\\[6pt]
  \rule{0.8\linewidth}{0.5pt}\\[6pt]
  {\large\bfseries Ancillary Botany-I}\\[4pt]
  {\large\bfseries Paper: BOB-203A}\\[4pt]
  {
\normalsize Time: Three Hours \hfill Maximum Marks: 70}
\end{center}

\rule{\linewidth}{0.4pt}

\smallskip



\noindent   \textit{Note: Attempt five questions in all, including question no. 1 which is compulsory. The figures in the right-hand margin indicate marks.}

\bigskip

\medskip


\noindent   \textbf{Question 1.} Fill in the blanks or answer the following in not more than 50 words each:\pts{10}

\begin{parts}
  \item Why is RNA considered as the progenitor for the origin of first cellular life on the earth?
  \item How old is the Earth?
  \item 'Three-domain' system of classification was proposed by \makebox[3cm]{\hrulefill}.
  \item Heterocyst is the site of \makebox[3cm]{\hrulefill}.
  \item Name two microbes used as food.
  \item Mention a typical feature to differentiate the Angiosperms from the Gymnosperms.
  \item What is idioandrosporous?
  \item What is the function of elater?
  \item What is a secondary pollutant?
  \item Define wet deposition.
\end{parts}

\medskip


\noindent   \textbf{Question 2.} Write short notes on any two of the following:\pts{15}

\begin{parts}
  \item Microorganisms capable of producing antibiotics
  \item Cyanobacteria as biofertilizer
  \item Paper / Bacterial endospores
\end{parts}

\medskip


\noindent   \textbf{Question 3.} Give a brief account of the classification of microbes as proposed by Whittaker.\pts{15}

\medskip


\noindent   \textbf{Question 4.} Briefly describe the structure and asexual reproduction of   \textit{Oedogonium}.\pts{15}

\medskip


\noindent   \textbf{Question 5.} Draw the labelled diagram of any two of the following (description not required):\pts{15}

\begin{parts}
  \item T.S. stem of   \textit{Selaginella}
  \item   \textit{Marchantia} sporophyte
  \item Thallus structure of   \textit{Rhizopus}
\end{parts}

\medskip


\noindent   \textbf{Question 6.} Briefly describe the reproduction in Gymnosperms with special reference to   \textit{Cycas}.\pts{15}

\medskip


\noindent   \textbf{Question 7.} Give an account of air pollutants and their sources. Suggest suitable methods for controlling air pollution.\pts{15}

\medskip


\noindent   \textbf{Question 8.} Define the sources of greenhouse gases in the atmosphere. Explain the greenhouse effect and its impact on climate change.\pts{15}

\medskip


\noindent   \textbf{Question 9.} Write short notes on any two of the following:\pts{15}

\begin{parts}
  \item V.S. archegoniophore of   \textit{Marchantia}
  \item Acid rain
  \item Threats to biodiversity
\end{parts}

\vfill
\rule{\linewidth}{0.4pt}

\begin{center}\small
\href{https://bhu-exam.vercel.app/}{Click Here}\\[4pt]\href{https://me-aryan.vercel.app/}{Click Here}\\[4pt]\href{https://quashan-me.vercel.app/}{Click Here}
\end{center}

\end{document}
